\section{The 3+3 resolvent}

Put
\[
 P_X(T)=T^6+2T^5+4T^4+6T^3+(X-5)T^2-8T-16
\]
so that $P_\zeta(T)=F_\zeta(T)$.  Let $r_1,\ldots,r_6$ be the roots of
$P_X$ in a splitting field over $\mathbf Q(X)$.  If
$A\subset\{1,\ldots,6\}$ has cardinality three,
write $y_A=\prod_{i\in A}r_i$.  For an unordered complementary pair
$\{A,A^c\}$, put $s_A=y_A+y_{A^c}$.  There are ten such pairs, and we
define
\[
 \mathcal R_X(S)=\prod_{\{A,A^c\}}(S-s_A).
\]
The coefficients are symmetric polynomials in the $r_i$ with integer
coefficients.
The fundamental theorem on symmetric polynomials therefore gives
$\mathcal R_X\in\mathbf Z[X,S]$.

Let $C_X$ be the companion matrix of $P_X$, in a convention for which
its characteristic polynomial is $P_X$, and put
\[
 Q_X(Y)=\det\!\left(YI-\bigwedge\nolimits^3C_X\right).
\]
The eigenvalues of $\bigwedge^3C_X$ are the twenty numbers $y_A$.
Moreover,
\[
 y_Ay_{A^c}=r_1\cdots r_6=-16,
 \qquad
 y_{A^c}=-\frac{16}{y_A}.
\]
Consequently, for each complementary pair,
\[
 (Y-y_A)(Y-y_{A^c})=Y^2-s_AY-16,
\]
and hence
\begin{equation}\label{eq:exterior-resolvent}
 Q_X(Y)=Y^{10}\mathcal R_X\!\left(Y-\frac{16}{Y}\right).
\end{equation}

Exact expansion gives
\begin{align}
\mathcal R_X(S)={}&S^{10}+6S^9+4(X+35)S^8
 +2(X^2-6X+389)S^7\notag\\
& +(X^3-15X^2+667X+5683)S^6
 +96(X^2-2X+337)S^5\notag\\
& +64(X^3-14X^2+477X+800)S^4
 -1536(X^2-21X-368)S^3\notag\\
& +8192(X^2+35X-48)S^2
 -65536(X-16)(X+4)S\notag\\
& -65536(X-16)^2.
\label{eq:explicit-resolvent}
\end{align}
Identity \eqref{eq:explicit-resolvent} is certified independently by
the exact-arithmetic exterior-power computation described in Appendix
\ref{sec:exact-certificates}; no numerical specialization is used.

\section{Absence of K-rational roots of the resolvent}

Let $\zeta\ne1$ be a root of unity, put $K=\mathbf Q(\zeta)$, and set
$\mathcal R_\zeta=\mathcal R_X|_{X=\zeta}$.  Fix a prime
$\mathfrak p\mid2$, normalize $v=v_{\mathfrak p}$ by
$v(K_{\mathfrak p}^\times)=\mathbf Z$, and put $e=v(2)$.
Write $\bar x$ for the reduction of an integral element $x$ modulo
$\mathfrak p$.

\subsection{Orders with a nontrivial odd part}

Write $\operatorname{ord}(\zeta)=2^am$ with $m>1$ odd, and decompose
$\zeta=\xi\eta$, where $\xi$ has $2$-power order and $\eta$ has odd
order.  Since $\bar\xi=1$ and reduction is injective on odd-order roots
of unity, $\bar\zeta=\bar\eta\ne1$.

Write $\mathcal R_\zeta(S)=\sum_{j=0}^{10}c_jS^j$.  Reducing the unit
part of each coefficient in \eqref{eq:explicit-resolvent} gives the
following exact valuations:
\[
 \bigl(v(c_0),\ldots,v(c_{10})\bigr)
 =e(16,16,13,9,6,5,0,1,2,1,0).
\]
For completeness, the reductions that prove exactness are
\begin{align*}
\left(c_j/2^{v(c_j)/e}\bmod\mathfrak p\right)_{j=0}^{10}
={}&\bigl(\bar\zeta^2,\bar\zeta^2,
 \bar\zeta(\bar\zeta+1),\bar\zeta(\bar\zeta+1),\\
&\quad \bar\zeta(\bar\zeta+1)^2,(\bar\zeta+1)^2,
 (\bar\zeta+1)^3,(\bar\zeta+1)^2,\\
&\quad \bar\zeta+1,1,1\bigr).
\end{align*}
Odd rational-integer factors reduce to $1$ and have therefore been
suppressed; every displayed residue is nonzero.

The lower convex hull is
\begin{equation}\label{eq:odd-res-np}
 (0,16e)\longrightarrow(6,0)\longrightarrow(10,0).
\end{equation}
Indeed, after division by $e$, the differences between the valuations at
abscissas $1,\ldots,5$ and the line from $(0,16e)$ to $(6,0)$ are
\[
 \frac83,\quad\frac73,\quad1,\quad\frac23,\quad\frac73,
\]
all positive.  From abscissa $6$ onward the lower boundary is horizontal
at height $0$.  The first six roots thus have
valuation $8e/3$.  The ramification index $e$ is a power of two, so
$8e/3\notin\mathbf Z$; none of these six roots belongs to $K_{\mathfrak p}$.

It remains to exclude the four unit roots.  Since $\gcd(4,m)=1$, the map
$x\mapsto x^4$ is an automorphism of the cyclic group $\langle\eta\rangle$.
Let $\theta\in\langle\eta\rangle$ be the unique element satisfying
$\theta^4=\eta$, and put
\[
 s_0=(\theta+1)^3.
\]
Since $\theta\ne1$ and reduction is injective on odd-order roots of unity,
$\bar\theta\ne1$; hence $s_0$ is a $\mathfrak p$-adic unit.  Only the coefficients
of $S^{10}$ and $S^6$ survive modulo $\mathfrak p$, and
\[
 \overline{c_6}=(\bar\zeta+1)^3
 = (\bar\theta+1)^{12}=\bar s_0^4.
\]
Thus
\begin{equation}\label{eq:odd-res-reduction}
 \overline{\mathcal R_\zeta(S)}
 =S^6(S-\bar s_0)^4.
\end{equation}

We shall use the following three coefficientwise congruences in
$\mathbf Z[u]$:
\begin{align}
 \mathcal R_{u^4}\bigl((u+1)^3\bigr)
 &\equiv2u^3(u+1)^{27}\pmod4,
 \label{eq:cong-I}\\
 \left.
 \frac{\partial\mathcal R_X}{\partial X}\bigl((u+1)^3\bigr)
 \right|_{X=u^4}
 &\equiv(u+1)^{26}\pmod2,
 \label{eq:cong-II}\\
 \mathcal R_{-u^4}\bigl((u+1)^3\bigr)
 &\equiv2u^3(u+1)^{26}\pmod4.
 \label{eq:cong-III}
\end{align}
They are verified by the exact-arithmetic certificate described in
Appendix \ref{sec:exact-certificates}, which checks every coefficient of
the three differences.

We now prove
\begin{equation}\label{eq:res-s0-value}
 v\bigl(\mathcal R_\zeta(s_0)\bigr)=1.
\end{equation}
If $a=0$, then $\zeta=\eta=\theta^4$ and $e=1$; congruence
\eqref{eq:cong-I} writes $\mathcal R_\zeta(s_0)$ as twice a unit plus an
element of $4\mathcal O_{K_{\mathfrak p}}$, and proves
\eqref{eq:res-s0-value}.  If $a=1$, then $\zeta=-\eta=-\theta^4$ and
$e=1$, so \eqref{eq:cong-III} gives the same conclusion in the same way.

Suppose $a\ge2$.  Set $\delta=\eta(\xi-1)$, so that
$\zeta=\eta+\delta$ and $v(\delta)=1$.  Taylor expansion in the
polynomial variable $X$ gives
\[
 \mathcal R_{\eta+\delta}(s_0)
 =\mathcal R_\eta(s_0)
  +\delta\,\partial_X\mathcal R_X(s_0)|_{X=\eta}
  +\delta^2A
\]
for some $A\in\mathcal O_{K_{\mathfrak p}}$; its integrality follows
because $\mathcal R_X(s_0)$ has coefficients in
$\mathcal O_{K_{\mathfrak p}}$.  Congruence
\eqref{eq:cong-I} gives
$v(\mathcal R_\eta(s_0))=e\ge2$, while
\eqref{eq:cong-II} says that the displayed derivative is a unit.  The
linear term therefore has valuation one and every other term has
valuation at least two, proving \eqref{eq:res-s0-value} without
cancellation.

After the translation $S=s_0+Y$, all coefficients remain integral, and equation
\eqref{eq:odd-res-reduction} shows that the coefficients of
$Y,Y^2,Y^3$ have positive valuation and that the coefficient of $Y^4$
is a unit.  Together with \eqref{eq:res-s0-value}, this puts the edge
\[
 (0,1)\longrightarrow(4,0)
\]
in the Newton polygon.  Its slope is $-1/4$ and its horizontal length is
four.  Lemma \ref{lem:newton-denominator} therefore makes the
corresponding degree-four factor irreducible.  In particular, none of
these four roots lies in $K_{\mathfrak p}$.  Together with
\eqref{eq:odd-res-np}, this proves that $\mathcal R_\zeta$ has no root in
$K$ whenever $m>1$.

\subsection{Pure powers of two}

First let $\zeta=-1$.  Substitution in
\eqref{eq:explicit-resolvent} gives
$c_j=2^{v(c_j)}u_j$, with the following odd unit parts $u_j$:
\[
\begin{array}{c|rrrrrrrrrrr}
j&0&1&2&3&4&5&6&7&8&9&10\\ \hline
u_j
&-289&51&-41&519&77&255&625&99&17&3&1,
\end{array}
\]
and therefore
\[
 \bigl(v(c_0),\ldots,v(c_{10})\bigr)
 =(16,16,14,10,8,7,3,3,3,1,0).
\]
The differences above the line from $(0,16)$ to $(6,3)$ at abscissas
$1,\ldots,5$ are
\[
 \frac{13}{6},\quad\frac73,\quad\frac12,\quad\frac23,
 \quad\frac{11}{6}.
\]
The points at abscissas $7,8,9$ also lie strictly above the line from
$(6,3)$ to $(10,0)$.  Hence the lower convex hull is
\[
 (0,16)\longrightarrow(6,3)\longrightarrow(10,0),
\]
with slopes $-13/6$ and $-3/4$.  The six roots belonging to the first
edge have valuation $13/6$, and the four belonging to the second have
valuation $3/4$.  Neither valuation is integral.  In particular,
$\mathcal R_{-1}$ has no root in $\mathbf Q_2$.

Now suppose that $\zeta$ has order $2^a$ with $a\ge2$.  Put
$\pi=\zeta-1$ and $e=v(2)=2^{a-1}$, so $v(\pi)=1$.  Substituting
$X=1+\pi$ into the parenthesized coefficient polynomials in
\eqref{eq:explicit-resolvent} gives
\begin{align*}
 X^2+35X-48&=\pi^2+37\pi-12,\\
 X^2-21X-368&=\pi^2-19\pi-388,\\
 X^3-14X^2+477X+800&=\pi^3-11\pi^2+452\pi+1264,\\
 X^2-2X+337&=\pi^2+336,\\
 X^3-15X^2+667X+5683&=\pi^3-12\pi^2+640\pi+6336,\\
 X^2-6X+389&=\pi^2-4\pi+384,\\
 X+35&=\pi+36.
\end{align*}
In addition, $X-16=\pi-15$ and $X+4=\pi+5$ are units.  Since $e\ge2$,
each displayed expression has a unique term of least valuation.  It
follows that
\begin{align}
 \bigl(v(c_0),\ldots,v(c_{10})\bigr)
={}&(16e,16e,13e+1,9e+1,6e+2,5e+2,\notag\\
&\hspace{34mm}3,e+2,2e+1,e,0).
\label{eq:pure-res-valuations}
\end{align}
For the first segment, the differences at abscissas $1,\ldots,5$ are
\[
 \frac{16e-3}{6},\quad\frac{7e}{3},\quad
 \frac{2e-1}{2},\quad\frac{2e}{3},\quad
 \frac{14e-3}{6},
\]
all positive.  At abscissas $7,8,9$, the valuations $e+2,2e+1,e$ are
strictly greater than the respective heights $9/4,3/2,3/4$ of the
second segment.  Thus the lower convex hull is
\[
 (0,16e)\longrightarrow(6,3)\longrightarrow(10,0),
\]
with slopes
\[
 -\frac{16e-3}{6},\qquad -\frac34.
\]
The integer $16e-3$ is odd, and modulo three it is congruent to $e$.
Since $e$ is a power of two, $3\nmid e$; hence
\[
 \gcd(16e-3,6)=1.
\]
The two slopes have reduced denominators six and four, respectively,
and the corresponding horizontal lengths are six and four.  Lemma
\ref{lem:newton-denominator} therefore shows that the two slope factors
are irreducible of degrees six and four.  In particular,
$\mathcal R_\zeta$ has no root in $K_{\mathfrak p}$.

Combining the odd-part and pure-power cases proves that
\begin{equation}\label{eq:res-no-K-root}
 \mathcal R_\zeta(S)\text{ has no root in }K
 \qquad(\zeta\ne1).
\end{equation}
